\documentclass[pdftex,12pt,letterpaper]{report}
\usepackage[pdftex]{graphicx}
\usepackage[us, 12hr]{datetime}
\usepackage{verbatim}
\usepackage{moreverb}
\usepackage{hyperref}
\let\verbatiminput=\verbatimtabinput %tabs are ignored in verbatim, this corrects for that
\def\verbatimtabsize{4\relax} % set tabs=4 (else my output goes off the screen)

\write18{bash bin/producedata.sh examples/bettse.conf /Users/bettse/Dropbox/Finances/current.lgr}
%\write18{python degree_tree.py} %run to produce example graph file

\usepackage{fullpage}


\newcommand{\HRule}{\rule{\linewidth}{0.5mm}}
\newcommand{\insertplot}[1]{\includegraphics[width=0.9\linewidth, keepaspectratio=true]{build/#1}\\[1cm]}

\begin{document}

\begin{titlepage}

\begin{center}


% Upper part of the page
\includegraphics[width=0.33\textwidth]{templates/bettse/logo}\\[1cm]    

\textsc{\LARGE My Ledger}\\[1.5cm]

\textsc{\Large Eric Betts}\\[0.5cm]


% Title
\HRule \\[0.4cm]
{ \huge \bfseries A brief summary of my finances}\\[0.4cm]

\HRule \\[1.5cm]

\vfill

% Bottom of the page
{\large \today }
{\large \currenttime }

\end{center}

\end{titlepage}


\tableofcontents

\chapter*{Introduction}

"I kept account of every farthing I spent, and my expenses were carefully calculated. 
Every little item, such as omnibus fares or postage or a couple of coppers spent on newspapers, would be entered, and the balance struck every evening before going to bed.
That habit has stayed with me ever since, and I know that as a result, though I have had to handle public funds amounting to lakhs, I have succeeded in exercising strict economy in their disbursement, and instead of outstanding debts have had invariably a surplus balance in respect of all the movements I have led.
Let every youth take a leaf out of my book and make it a point to account for everything that comes into and goes out of his pocket, and like me he is sure to be a gainer in the end.
" -- M.K.Gandhi autobiography

\chapter{Budget}

\insertplot{month_breakdown}

\verbatiminput{/Users/bettse/Dropbox/Finances/budget.lgr}

\section{Current month progress}

Negative values indicate to-be-spent funds.  Positive vaues indicate overspending.

\verbatiminput{build/budget.txt}

\chapter{Assets}

\verbatiminput{build/assets.txt}

\chapter{Liabilities}

\verbatiminput{build/liabilities.txt}

\chapter{Transactions}

\section{First Tech Credit Union}

Balance over the last year to date :

\insertplot{checking-1yearbalance}

Balance during \monthname :

\insertplot{checking-1monthbalance}

Transactions for the last 7 days:

\verbatiminput{build/checking-1trans.txt}

\chapter{Net Worth}

\begin{itemize}

\item The balance of my assets to my liabilities gives my net worth (including retirement funds): \input{build/networth.txt}
\item Removing long term investment and loan accounts gives my net liquidity: \input{build/liquidity.txt}
\item Balancing expenses against income yields your cash flow, or net profit/loss(negative is profit, positive is loss): \input{build/cashflow.txt}

\end{itemize}

\chapter{Forecast}

\section{First Tech Credit Union}

\insertplot{checking-1forecast}

\verbatiminput{build/checking-1forecast.txt}

\end{document}
